\section{Related Work}
\label{sec:related_work}

Our proposed framework lies at the intersection of parameter-space model merging, online test-time adaptation, and second-order loss landscape analysis. In this section, we review the mathematical foundations of these fields and highlight where our work departs from and improves upon existing literature.

\subsection{Parameter-Space Model Merging}
Model merging has emerged as a prominent paradigm for combining task-specific capabilities of neural networks without invoking additional training costs. Foundational works such as Model Soups \citep{wortsman2022model} and Task Arithmetic \citep{ilharco2022editing} established that fine-tuning a shared pre-trained base model $\theta_0$ on different tasks creates weight configurations that lie within a shared, low-loss basin. Consequently, task-specific task vectors $v_k = \theta_k - \theta_0$ can be linearly interpolated. 

However, scaling to multiple tasks often introduces inter-task representation interference, where parameter-wise conflicts (e.g., sign disagreements) degrade joint performance. To resolve this, TIES-Merging \citep{yadav2023ties} introduces a heuristic sign-agreement and thresholding protocol, while DARE \citep{yu2023language} uses random delta-dropout to sparsify task vectors before merging. Other works like RegCalMerge \citep{regcalmerge} calibrate model outputs via class-capacity normalization. 

While effective, these approaches apply static, uniform scaling factors across the entire network architecture. To introduce depth-dependent specialized scaling, PolyMerge \citep{polymerge} projects the merging coefficients into a rigid low-dimensional continuous polynomial subspace. In contrast, RCR-Merge allows full layer-wise adaptability but bounds the spatial trajectory of coefficients using the pre-trained base model's second-order geometry, striking a superior balance between adaptability and representational smoothness.

\subsection{Adaptive and Test-Time Model Merging}
Rather than relying on static, pre-computed merging weights, \emph{Adaptive Model Merging} frameworks optimize merging coefficients post-hoc. \citet{adamerging} introduced AdaMerging, which parameterizes the merging process with layer-wise coefficients $\boldsymbol{\lambda} \in \mathbb{R}^{K \times L}$. During deployment, AdaMerging adapts these coefficients using unsupervised test-time adaptation (TTA). The primary optimization objective is the minimization of Shannon entropy of predicted probability distributions, a technique pioneered by Tent \citep{wang2020tent} for source-free domain adaptation \citep{liang2023source}.

However, online TTA is notorious for its vulnerability to transductive noise, small batch sizes, and distribution shifts, which often lead to degenerate constant-prediction states or catastrophic representation collapse \citep{liao2021are}. While traditional TTA literature introduces temporal model-averaging or heavy normalization constraints to stabilize predictions, the specific challenges of adaptive model merging—where we optimize a low-dimensional coefficient space $\boldsymbol{\lambda}$ that directly scales high-dimensional parameter updates—have been ignored. We are the first to identify and formalize the Overfitting-Optimizer Paradox in this context, demonstrating that unconstrained entropy minimization in the coefficient space fits transductive noise, and we provide the first second-order geometric solution to this paradox.

\subsection{Loss Landscape Geometry and Fisher Information}
Our geometric regularizer is deeply grounded in the second-order mathematical structure of neural loss landscapes. Foundational work in generalization theory demonstrates that flatter minima—characterized by low eigenvalues of the loss Hessian—generalize significantly better than sharp minima under distribution shifts \citep{foret2020sharpness, cha2021swad, izmailov2018averaging}. 

To measure local curvature without performing expensive Hessian-vector products, literature often relies on the Fisher Information Matrix (FIM), which acts as a localized Riemannian metric on the parameter manifold. \citet{matena2021merging} first leveraged the diagonal FIM of fine-tuned models to perform Fisher-weighted averaging, establishing that parameter-wise distance should be scaled by local curvature. While their work focuses on offline, static merging, we pivot this philosophy to online, adaptive test-time optimization. By using the pre-trained base model's FIM trace as a spatial Riemannian metric across depth, RCR-Merge is the first to demonstrate that pre-trained base model curvature can be used to construct a highly stable and robust analytical barrier for test-time adaptation.